\documentclass[11pt,reqno]{amsart}
%=====================================================================
%  TEAM A1 — INTEGRATED WORKING PAPER (Week 1 state, ~70–75% to MS bar)
%  Flexible Computing Loads, Grid Resilience, and Carbon Emissions.
%  Weaves together all four A1 roles:
%   R1 General Carbon Theory · R2 Solver/Numerical · R3 Comparative
%   Statics · R4 Theory Hardening & Regimes.
%  Honest about what is proved vs. conditional/conjectural.
%=====================================================================
\usepackage[T1]{fontenc}\usepackage{lmodern}
\usepackage{amsmath,amssymb,amsthm,mathtools}
\usepackage[margin=1in]{geometry}
\usepackage{microtype}\usepackage{enumitem}\usepackage{booktabs}\usepackage{xcolor}
\usepackage[colorlinks=true,linkcolor=blue!55!black,citecolor=blue!55!black]{hyperref}

\theoremstyle{plain}
\newtheorem{theorem}{Theorem}\newtheorem{proposition}{Proposition}
\newtheorem{lemma}{Lemma}\newtheorem{corollary}{Corollary}
\theoremstyle{definition}\newtheorem{condition}{Condition}\newtheorem{assumption}{Assumption}
\theoremstyle{remark}\newtheorem{remark}{Remark}

\newcommand{\KR}{K_R}\newcommand{\KF}{K_F}\newcommand{\KM}{K_M}
\newcommand{\KRs}{K_R^{*}}\newcommand{\KFs}{K_F^{*}}\newcommand{\KMs}{K_M^{*}}
\newcommand{\KRo}{K_R^{\circ}}\newcommand{\KFo}{K_F^{\circ}}
\newcommand{\Dbar}{\bar D}\newcommand{\CM}{C_M}\newcommand{\CF}{C_F}\newcommand{\CR}{C_R}
\newcommand{\CMbar}{\bar C_M}\newcommand{\bp}{\beta+p}\newcommand{\FKR}{F_{\!\KR}}
\newcommand{\E}{\mathbb{E}}\newcommand{\PP}{\mathbb{P}}\newcommand{\eqdef}{\vcentcolon=}
\newcommand{\Ucal}{\mathcal U}

\begin{document}

\title[Flexible Computing Loads, Grid Resilience, and Carbon Emissions]
{Flexible Computing Loads, Electricity-Grid Resilience, and Carbon Emissions:\\
A General Theory of When Mining and AI Training Reduce Emissions}
\author{Team A1 (Workstream A: Optimal Contracting Theory and Quantitative Validation)}
\date{Integrated working draft}

\begin{abstract}
Intensive, interruptible computing demands---Bitcoin mining and, imperfectly, AI
training---can be dispatched in real time. We ask when introducing such a
\emph{flexible computing load} (\textsc{fcl}) into a regulated grid \emph{reduces}
carbon emissions. The program changes the utility's optimal mix of renewable
($\KR$), fossil ($\KF$), and computing ($\KM$) capacity, setting up two opposing
forces: a distribution-free \emph{utilization effect} (computing burns spare
fossil, raising emissions) and a \emph{capacity-substitution effect} (computing
revenue $R>c$ raises the shadow value of renewables, so $\KRs>\KRo$, lowering
emissions). Our contributions: (i) we prove the capacity-substitution channel
$\KRs>\KRo$ for \emph{any} residual-demand distribution satisfying a transparent
shift-ratio condition plus interior regularity; (ii) we give the exact net-carbon
inequality and a clean sufficient condition; (iii) we identify a closed-form
\emph{viability threshold} on the renewable cost, with a rigorous general
large-cost anchor and an explicit power-law window; and (iv) at fixed equilibrium
capacities we reduce the carbon effect to a two-mean affine form, bracket it by
closed-form bounds, and identify the uniform as the least-favorable conditional
shape. Every analytical formula is verified to machine precision, and the
distribution-free claims are confirmed across eight distribution families. We are
explicit about what remains conditional: the fully general, finite-optimum carbon
guarantee is proved for the power-law family and conjectured beyond it.
\end{abstract}

\maketitle

%---------------------------------------------------------------------
\section{Introduction}\label{sec:intro}

Bitcoin mining consumes electricity on a national scale, and AI training is on a
similar trajectory. Both are unusual loads: they can be ramped down in seconds at
little cost, so a grid operator can treat them as a real-time \emph{demand-response}
resource. A provocative possibility follows---that adding such a load could
\emph{lower} carbon emissions, by making renewable investment more profitable (the
computing load monetizes otherwise-curtailed renewable surplus) even as it also
burns some fossil fuel when it runs on spare thermal capacity. Whether the net
effect is lower or higher emissions is the question of this paper.

We study a regulated utility that chooses renewable capacity $\KR$, dispatchable
fossil backup $\KF$, and flexible-computing-load capacity $\KM$ to maximize
expected profit, and dispatches the computing load by an optimal real-time rule.
Introducing the program sets up two forces on carbon. The \textbf{utilization
effect} is that computing run on spare fossil capacity strictly raises fossil
output---this always pushes emissions up and is distribution-free. The
\textbf{capacity-substitution effect} is that the computing margin $R-c>0$ raises
the shadow value of renewables, so the utility builds \emph{more} renewables
($\KRs>\KRo$), which shifts residual demand toward surplus and pushes emissions
down. The carbon question is which force wins.

\paragraph{Contributions and roadmap.} This integrated draft is the joint output
of four roles. \S\ref{sec:model} sets up the model and the optimal dispatch
contract. \S\ref{sec:channel} proves the capacity-substitution channel
$\KRs>\KRo$ in general (Role~1). \S\ref{sec:carbon} gives the exact net-carbon
condition and its distribution-free reading (Role~1). \S\ref{sec:threshold}
establishes the renewable-cost viability threshold (Role~1).
\S\ref{sec:regimes} maps the regimes, completes the small-\textsc{fcl} case, and
pressure-tests the maintained assumptions (Role~4). \S\ref{sec:numerics} reports
the solver, the machine-precision audit, and the comparative statics across eight
distribution families (Roles~2--3). \S\ref{sec:open} states honestly what is
proved versus conditional, and the path to a complete theory.

\paragraph{Scope.} This is the \emph{general theory and numerical methods}
contribution. Empirical calibration to real grids, the reliability/welfare and
social-versus-private analysis, and the generalization to a flexibility continuum
(storage, partially-flexible AI, multi-grid markets, mining dynamics) are
companion efforts and are out of scope here.

%---------------------------------------------------------------------
\section{Model and the optimal dispatch contract}\label{sec:model}

A regulated utility serves a representative period. The state is \emph{residual
demand}
\begin{equation}\label{eq:resdem}
  \tilde x \eqdef \tilde D - \KR\tilde S \in[-\KR,\Dbar],
\end{equation}
where $\tilde D\in[0,\Dbar]$ is stochastic consumer demand and $\tilde S\in[0,1]$
the renewable capacity factor; $F(x;\KR)$ is the CDF with density $f$, extended by
$F\equiv0$ below $-\KR$ and $F\equiv1$ above $\Dbar$. More renewables shift mass
left: $\FKR\eqdef\partial F/\partial\KR\ge0$. Per-MWh: computing revenue $R$,
fossil cost $c$, retail price $p$, blackout penalty $\beta$. Amortized per-MW
capital: $\CR,\CF,\CM$.

\begin{assumption}[Regularity; maintained throughout]\label{ass:reg}
$R>c+\CM$,\quad $R<c+\CF+\CM$,\quad $\bp>c+\CF$,\quad $\bp>R$.
\end{assumption}

\noindent A \emph{demand-response contract} is a dispatch policy
$d:[-\KR,\Dbar]\to[0,\KM]$. Computing can only run on power physically available
(spare fossil plus renewable surplus), $a(x)\eqdef(\KF-x)_+$.

\begin{lemma}[Optimal contract; Role~1]\label{lem:dispatch}
For any $F$ and any capacities, the unique optimal served dispatch is the
renewable-first threshold rule
$d^{*}(x)=\min\{(\KF-x)_+,\KM\}$: full computing on deep surplus, exact
gap-filling on $(\KF-\KM,\KF)$, zero in scarcity.
\end{lemma}

Substituting $d^{*}$ and integrating by parts gives the profit
\begin{equation}\label{eq:profit}
  \Pi=p\E[\tilde D]+R\!\!\int_{\KF-\KM}^{\KF}\!\!F\,dx
  -c\!\!\int_{-\KM}^{\KF-\KM}\!\!(1-F)\,dx
  -(\bp)\!\!\int_{\KF}^{\Dbar}\!\!(1-F)\,dx-\CF\KF-\CR\KR-\CM\KM,
\end{equation}
with first-order conditions (valid in both regimes; Roles~1,4)
\begin{align}
  \tfrac{\partial\Pi}{\partial\KM}&=(R-c)F(\KF-\KM)+cF(-\KM)-\CM,\label{eq:gM}\\
  \tfrac{\partial\Pi}{\partial\KF}&=R[F(\KF)-F(\KF-\KM)]-c[1-F(\KF-\KM)]+(\bp)[1-F(\KF)]-\CF,\label{eq:gF}\\
  \tfrac{\partial\Pi}{\partial\KR}&=R\!\!\int_{\KF-\KM}^{\KF}\!\!\FKR+c\!\!\int_{-\KM}^{\KF-\KM}\!\!\FKR+(\bp)\!\!\int_{\KF}^{\Dbar}\!\!\FKR-\CR.\label{eq:gR}
\end{align}

\begin{lemma}[Distribution-free cutoffs; Roles~1,4]\label{lem:cutoffs}
The no-\textsc{fcl} threshold satisfies $F(\KFo;\KRo)=\varphi^{\circ}\eqdef(\bp-c-\CF)/(\bp-c)$,
independent of $F$. In the large-\textsc{fcl} regime ($F(-\KMs)=0$), also
$F(\KFs;\KRs)=(\bp-c-\CF-\CM)/(\bp-R)$ and $F(\KFs-\KMs;\KRs)=\CM/(R-c)$. Hence
the adoption threshold $\CMbar=(R-c)\varphi^{\circ}$.
\end{lemma}

%---------------------------------------------------------------------
\section{The capacity-substitution channel}\label{sec:channel}

The carbon result rests on $\KRs>\KRo$: the program builds more renewables.
The driver is visible in \eqref{eq:gR}: on the dispatch band $(\KF-\KM,\KF)$ the
shadow value of renewables is weighted by $R$, whereas without computing it is
weighted by $c<R$.

\begin{condition}[Shift ratio]\label{cond:cs}
$g(x)\eqdef\FKR(x)/f(x)=\E[\tilde S\mid\tilde x=x]$ is non-increasing on the
dispatch window at each point of the adoption path.
\end{condition}

\begin{condition}[Regularity]\label{cond:reg}
Along the adoption path the program has an interior optimum at which the Hessian
$\mathcal H$ of $\Pi$ is nonsingular (hence $\prec0$).
\end{condition}

\begin{theorem}[Capacity-substitution channel; Role~1]\label{thm:channel}
Under Conditions~\ref{cond:cs}--\ref{cond:reg}, renewable capacity is strictly
higher with the program: $\KRs>\KRo$ whenever $\CM<\CMbar$. The mechanism is the
comparative static
$\;d\KRs/d\CM = (R-\bp)(R-c)\,[\FKR(\KFs)f(\KFs-\KMs)-\FKR(\KFs-\KMs)f(\KFs)]/\det\mathcal H\le0$,
which integrates from $\CMbar$ (where $\KMs=0$, $\KRs=\KRo$) downward.
\end{theorem}

The proof differentiates the FOC system implicitly; the cofactor sign reduces to
$g(\KFs)-g(\KFs-\KMs)\le0$, exactly Condition~\ref{cond:cs}. The obstruction to a
one-line lattice argument is that $\Pi$ is \emph{submodular} in $(\KR,\KF)$ while
having increasing differences in $\KM$; the homotopy sidesteps it.

\begin{corollary}[Power-law family; explains footnote-6]\label{cor:power}
For $F(x;\KR)=((x+\KR)/(\Dbar+\KR))^{\alpha}$, $g(x)=(\Dbar-x)/(\Dbar+\KR)$ is
affine and decreasing for \emph{every} $\alpha>0$, so $\KRs>\KRo$ holds throughout
the family.
\end{corollary}

\begin{remark}[Reading; sharpness]
$g$ is the conditional renewable factor, so Condition~\ref{cond:cs} says
high-stress states are renewable-poor---a negative demand--supply association. The
cofactor bracket flips sign with $g$, so an increasing $g$ would reverse the
channel; constructing a \emph{realized} reversed-channel example is left to the
robustness round (Role~3).
\end{remark}

%---------------------------------------------------------------------
\section{The net-carbon condition}\label{sec:carbon}

Carbon equals $c$ times expected fossil output. Integrating \eqref{eq:profit}'s
fossil terms by parts (Role~1) gives the exact identities
$\E[\text{fossil}^{*}]=\KFs-\tfrac{\CM}{R-c}(\KFs-\KMs-\mu^{*}_{\mathrm{low}})$ and
$\E[\text{fossil}^{\circ}]=\KFo(1-\varphi^{\circ})+(\varphi^{\circ}-q^{\circ})\mu^{\circ}_{\mathrm{mid}}$,
with $q^{\circ}=F(0;\KRo)$ and conditional means $\mu^{*}_{\mathrm{low}},\mu^{\circ}_{\mathrm{mid}}$.

\begin{proposition}[Exact net-carbon inequality; Role~1]\label{prop:C}
Carbon falls with the program iff
\begin{equation}\tag{C}\label{eq:C}
  \underbrace{\tfrac{\CM}{R-c}(\KFs-\KMs-\mu^{*}_{\mathrm{low}})}_{\text{capacity-substitution}}
  >
  \underbrace{\KFs-\KFo(1-\varphi^{\circ})-(\varphi^{\circ}-q^{\circ})\mu^{\circ}_{\mathrm{mid}}}_{\text{net utilization}}.
\end{equation}
The utilization effect is always positive and distribution-free; the
capacity-substitution effect is signed by Theorem~\ref{thm:channel}.
\end{proposition}

\begin{corollary}[Sufficient conditions; Role~1]\label{cor:suff}
If $f(\cdot;\KRs)$ is non-increasing on the \textsc{fcl} zone and $f(\cdot;\KRo)$
non-decreasing on the no-\textsc{fcl} fossil zone (shape class $\Ucal$; the uniform
is the boundary), then (C) holds at least as strongly as under the uniform.
\end{corollary}

\paragraph{A distribution-free reading (at fixed capacities; Role~1).} Holding the
equilibrium capacities and $q^{\circ}$ fixed, $\Delta\eqdef\E[\text{fossil}^{\circ}]-\E[\text{fossil}^{*}]$
is \emph{affine} in $(\mu^{*}_{\mathrm{low}},\mu^{\circ}_{\mathrm{mid}})$, hence
bracketed by closed-form bounds
$\KFo(1-\varphi^{\circ})-\KFs\le\Delta\le\KFo-\KFs+\tfrac{\CM}{R-c}(\KFs-\KMs+\KRs)$
computable from cutoff data; and within $\Ucal$ the uniform minimizes $\Delta$
(it maximizes $\mu^{*}_{\mathrm{low}}$ and minimizes $\mu^{\circ}_{\mathrm{mid}}$).
These are fixed-capacity statements; turning the last into a genuine cross-$F$
certificate requires signing the capacity covariation, which we leave open
(\S\ref{sec:open}).

%---------------------------------------------------------------------
\section{The renewable-cost viability threshold}\label{sec:threshold}

Write $\Delta(\CR)$ for the net fossil saving at the optimum given renewable cost
$\CR$. Whether carbon falls is, in the end, about how cheap renewables are.

\begin{theorem}[Viability threshold; Role~1]\label{thm:threshold}
\textup{(i, general)} For all $\CR$ above a finite $B_0$, no renewables are built,
and the program differs from the no-\textsc{fcl} grid only by the (positive)
utilization effect, so $\Delta(\CR)=-U_0<0$: the carbon-reduction set is bounded
above for every $F$. \textup{(ii, power-law / knife-edge)} The thresholds are
explicit,
$\CR^{\dagger}=\tfrac12[c+\tfrac{(R-\CM-c)^2}{R-c}+\tfrac{(R-c-\CF-\CM)^2}{\bp-R}]$
and $\CR^{\circ}=\CF^2/2(\bp-c)$ with $\CR^{\circ}<\CR^{\dagger}$; on
$(\CR^{\circ},\CR^{\dagger})$ renewables are unbounded, $\E[\text{fossil}^{*}]\to0$,
and carbon strictly falls. Hence $\CR^{\dagger}=\sup\{\CR:\Delta\ge0\}$ is finite,
and positive for the power-law family.
\end{theorem}

The large-cost anchor (i) is fully general; the strict-fall window (ii) uses the
knife-edge (constant renewable marginal value) of the power-law family. For a
generic $F$ with an interior renewable optimum, the set $\{\Delta\ge0\}$ is bounded
above by (i), but its \emph{non-emptiness}---that carbon falls somewhere---is what
Condition~\ref{cond:cs} delivers in the cases we can sign and is conjectured in
general. Closing this is the keystone of the next round.

%---------------------------------------------------------------------
\section{Regimes and robustness of the assumptions}\label{sec:regimes}
\emph{(Role~4: theory hardening.)}

\paragraph{Small-\textsc{fcl} regime.} When $\KM<\KR$ the surplus term $cF(-\KM)$
in \eqref{eq:gM} does not vanish; the first computing unit is worth more, widening
adoption to
$\CMbar^{b}=\CMbar+c\,q^{\circ}>\CMbar$ (paper eq.~62), verified exactly.

\paragraph{Regime map.} The $\CM$-thresholds are strictly ordered,
\[
  \underbrace{R-c-\CF}_{\CM^{\mathrm{fc}}}<\underbrace{(R-c)\varphi^{\circ}}_{\CMbar}<\underbrace{\CMbar+c\,q^{\circ}}_{\CMbar^{b}},
\]
sorting parameter space into no adoption, small-\textsc{fcl}, large-\textsc{fcl}
interior, and full coverage; the boundedness axis $\CR^{\circ}<\CR^{\dagger}$ of
\S\ref{sec:threshold} composes with it into one statement.

\paragraph{Necessity of Assumption~\ref{ass:reg}.} $R>c+\CM$ guarantees an
interior $\KM$ ($\CM/(R-c)<1$); $\bp>R$ keeps the blackout cutoff well-defined and
gives consumers priority; $\bp>c+\CF$ makes fossil backup worth building. The
remaining bound $R<c+\CF+\CM$ is needed only for the \emph{interior} regime---
relaxing it lands in the well-defined full-coverage corner, not a contradiction.

%---------------------------------------------------------------------
\section{Numerical methods and comparative statics}\label{sec:numerics}
\emph{(Role~2: solver \& validation; Role~3: comparative statics.)}

\paragraph{Solver and exactness audit (Role~2).} An independent solver quadratures
\eqref{eq:profit} and solves \eqref{eq:gM}--\eqref{eq:gR} with no closed forms
assumed. It reproduces the uniform closed forms to machine precision and confirms
the draft's algebra, including three corrections (a dropped factor $c$ in the
printed carbon equations, a large-\textsc{fcl} flat term, and that the printed
knife-edge bracket equals $2\,\partial\Pi/\partial\KR$). An adversarial audit
checks all 15 analytical formulas (gradients, cross-partials, cofactor,
distribution-free cutoffs, fossil decomposition) against finite differences:
\textbf{15/15} agree to $6$--$16$ digits; the test suite passes \textbf{36/36}.

\paragraph{Comparative statics across distributions (Role~3).} Across eight
families (uniform; Beta$(0.5,3),(2,5),(2,2),(5,2)$; lognormal; exponential;
Pareto):
\begin{itemize}[leftmargin=1.5em,topsep=2pt,itemsep=1pt]
  \item the distribution-free profit and blackout halves of the adoption result
        hold \textbf{16/16} whenever $\CM<\CMbar$;
  \item the capacity-substitution channel $\KRs>\KRo$ holds \textbf{32/32};
  \item the utilization effect is positive in every scenario;
  \item with capacities frozen (the short-run/timing case), carbon weakly rises
        \textbf{96/96}, confirming the sign of that comparative static;
  \item carbon \emph{strictly falls} in \textbf{2/32} of the general-model
        scenarios, concentrated at low $\CR$ and right-skewed $F$---exactly the
        regime Theorem~\ref{thm:threshold} predicts.
\end{itemize}
The early reading ``carbon never falls'' and the theory ``carbon can fall'' are
reconciled (Role~4): the sweep sat above the buildout crossover where the
utilization effect dominates; the reduction obtains below $\CR^{\dagger}$ where
renewables are unbounded. The thin ``2/32'' is consistent with the result and is
the empirical leg the next round will strengthen into one curated, finite-optimum
narrative.

%---------------------------------------------------------------------
\section{Scope of the contribution and open problems}\label{sec:open}

What is \emph{proved}: the optimal dispatch contract (Lemma~\ref{lem:dispatch});
the distribution-free cutoffs (Lemma~\ref{lem:cutoffs}); the capacity-substitution
channel $\KRs>\KRo$ under Conditions~\ref{cond:cs}--\ref{cond:reg}
(Theorem~\ref{thm:channel}); the exact net-carbon inequality (C)
(Proposition~\ref{prop:C}); the fixed-capacity distribution-free reading; the
general large-cost anchor and the explicit power-law carbon window
(Theorem~\ref{thm:threshold}); the regime map and the necessity pattern of
Assumption~\ref{ass:reg}. All are machine-verified and cross-checked across roles.

What remains \emph{conditional or open}, and is the agenda for completion:
\begin{enumerate}[leftmargin=1.7em,topsep=2pt,itemsep=1.5pt]
  \item \textbf{General-$F$, finite-optimum carbon theorem (keystone).} Prove that
        $\{\Delta\ge0\}$ is non-empty for general $F$ under a \emph{primitive}
        shift-ratio condition plus a tail condition, with the strict fall occurring
        at a finite interior $\KRs$ (not the unbounded corner).
  \item \textbf{Existence/uniqueness and regularity from primitives.} Replace
        Condition~\ref{cond:reg}'s numerically-checked $\mathcal H\prec0$ with a
        primitive condition on $F$ (e.g.\ log-concavity) guaranteeing a unique
        interior optimum.
  \item \textbf{Cross-$F$ robust certificate.} Sign the capacity covariation so
        the least-favorable-distribution statement holds across $\Ucal$, not only
        at fixed capacities.
  \item \textbf{Positioning.} A related-work section on the methods used (monotone
        comparative statics, interruptible/demand-response economics,
        distributionally-robust analysis) and a bibliography.
\end{enumerate}
The spine is in place and internally consistent; these four items take the paper
from a strong working draft to a complete \emph{Management Science} submission.

\paragraph{Contributions by role.} Role~1 (general carbon theory): the contract,
the channel, condition (C), the threshold, the distribution-free reading. Role~2
(solver \& numerical validation): the independent solver, the 15/15 audit, the
corrections. Role~3 (comparative statics): the eight-family confirmations and the
carbon-sign map. Role~4 (theory hardening \& regimes): the small-\textsc{fcl}
case, the regime map, the assumption pressure-test, and the theory--numerics
reconciliation.

\end{document}
